\documentclass{article}
\usepackage{graphicx}
\usepackage{enumitem}
\usepackage{xcolor}
\usepackage{pifont} % Pour les symboles (croix)
\usepackage{pseudocode}
\usepackage{listings}
\usepackage[colorlinks=true,linkcolor=blue,urlcolor=blue,citecolor=blue]{hyperref}

\lstset{
    language=C,
    basicstyle=\ttfamily,
    keywordstyle=\color{blue},
    commentstyle=\color{green!60!black},
    stringstyle=\color{red},
    numbers=left,
    numberstyle=\tiny\color{gray},
    breaklines=true,
    showstringspaces=false,
    frame=single,
    tabsize=2
}

\begin{document}

\title{Rapport de projet ALGO}

{\centering
\author{Achraf,Thomas,Elouen,Woody}
\par}
\date{Janvier 2024}
\maketitle

%Si jamais vous voulez insérer une image prenez le modèle et modifiez le
%\begin{figure}
%  \centering
%  \includegraphics[width=0.8\linewidth]{imagze.jpg}
%  \caption{Image}
%  \label{fig:image1}
%\end{figure}

\thispagestyle{empty}

\newpage

\tableofcontents

\newpage

\maketitle

\section{\textbf{Introduction}}

\begin{itemize}[label={\textbullet}]
  \item Objectif 1 :
\end{itemize}

\newpage


\section{\textbf{Analyses}}

\subsection{Compression}

\subsection{Decompression}

\documentclass{article}
\usepackage{pseudocode}
\begin{document}

\section{TAD Code Binaire}
\begin{tad}
     \tadNom{Code Binaire}
        
     \tadDependances{TableDeCodage, Caractere }
     \begin{tadOperations}{obtenirCodeBinaire}
        \tadOperation{obtenirCodeBinaire}{\tadDeuxParams{TableDeCodage}{Caractere}}{\tadParam{CodeBinaire}}
        \tadOperationAvecPreconditions{estBinaire}{\tadParam{CodeBinaire}}{\tadParam{Booléen}}
     \end{tadOperations}
      \begin{tadSemantiques}{}
            obtenirCodeBinaire : Obtient grâce à la table de codage le code binaire d'un caractère 
    estBinaire : Permet de vérifier que le code est bien binaire
         
      \end{tadSemantiques}
    \begin{tadPreconditions}{obtenirFilsGauche(A)}
            \tadPrecondition{estBinaire(code)}{non estVide(code)}
        \end{tadPreconditions}

\end{tad}   


\section{TAD Octet}
\begin{tad}
\tadNom{Octet}
\tadDependances{Chaine de caracteres, Code Binaire, Booleen}
\begin{tadOperations}{estOctet}
    \tadOperation{creerOctet}{}{\tadParams{Octet}}
    \tadOperation{ajouter}{\tadDeuxParams{Octet}{Bit}}{\tadParams{Octet}}
    \tadOperationAvecPreconditions{supprimer}{\tadDeuxParams{Octet}{Bit}}{\tadParams{Octet}}
    \tadOperation{estOctet}{\tadParams{Octet}}{\tadParams{Booleen}}
    \tadOperation{longueur}{\tadParams{Octet}}{\tadParams{\naturel}}
    \tadOperationAvecPreconditions{encoder}{\tadDeuxParams{Octet}{Code Binaire}}{\tadParams{ Chaine de caracteres}}
    \tadOperationAvecPreconditions{decoder}{\tadDeuxParams{Chaine de caracteres}{Code Binaire}}{\tadParams{Octet}}
    \tadOperationAvecPreconditions{obtenirOctet}{\tadParams{Chaine de caracteres}}{\tadParams{Octet}}
\end{tadOperations}
    
\end{tad}
    
\section{TAD Statistiques}
\begin{tad}
\tadNom{Statistiques}
\tadDependances{Caractère, Chaine de caracteres, Naturel}

\begin{tadOperations}{statistiques}

\tadOperation{statistiques}{\tadDeuxParams{{\caractere}{\naturel}}}{\tadParams{ Statistiques}}
\tadOperation{obtenirOccurences}{\tadParams{Statistiques}}{\tadParams{\naturel}}
\tadOperation{obtenirCaractere}{\tadParams{Statistiques}}{\tadParams{\caractere}}
\end{tadOperations}
\begin{tadSemantiques}{}
     statistiques : créer une statistique à partir d'un caractère et de son nombre d'occurences. \newline
     {obtenirOccurences} : Renvoie le nombre d'occurences d'un caractère et le stock dans la statistique    \newline
\end{tadSemantiques}

                          
\end{tad}



\section{TAD ArbreHuffman}

    \begin{tad}
        \tadNom{Arbre Huffman}
        
        \tadDependances{statistiques}
        
        \begin{tadOperations}{obtenirPonderation}
        
            \tadOperation{arbrehuffman}{Statistiques}{\tadParams{Arbre Huffman}}
            
            \tadOperationAvecPreconditions{fusionner}{\tadDeuxParams{Arbre Huffman}{Arbre Huffman}}{\tadParam{Arbre Huffman}}

            \tadOperation{obtenirPonderation}{\tadParams{Arbre Huffman}}{\tadParams{\naturel}}

            \tadOperationAvecPreconditions{fixerFilsDroit}{\tadDeuxParams{\paramEntree{arbred : Arbre Huffman}\paramEntreeSortie{racine : Arbre Huffman}}}

            \tadOperationAvecPreconditions{fixerFilsGauche}{\tadDeuxParams{\paramEntree{arbreg : Arbre Huffman}\paramEntreeSortie{racine : Arbre Huffman}}}

            \tadOperationAvecPreconditions{obtenirFilsGauche}{\tadParams{Arbre Huffman}}{\tadParams{Arbre Huffman}}

            \tadOperationAvecPreconditions{oHbtenirFilsDroit}{\tadParams{Arbre Huffman}}{\tadParams{Arbre Huffman}}

            \tadOperation{estFeuille}{\tadParams{Arbre Huffman}}{\tadParams{\booleen}}

            \tadOperation{estVide}{\tadParams{Arbre Huffman}}{\tadParams{\booleen}}
        \end{tadOperations}
        
        \begin{tadSemantiques}{}
        
            
            
        \end{tadSemantiques}
        
        \begin{tadPreconditions}{obtenirFilsGauche(A)}
            \tadPrecondition{fusionner(A)}{estVide(A)}
            \tadPrecondition{fixerFilsDroit(r,ad)}{non(estVide(r))}
            \tadPrecondition{fixerFilsGauche(r,ag)}{non(estVide(r))}
            \tadPrecondition{obtenirFilsGauche(A)}{non estUneFeuille(A)}
            \tadPrecondition{obtenirFilsDroit(A)}{non estUneFeuille(A)}
        \end{tadPreconditions}
           
        \end{tad}

    


\section{FileDePriorité}
\begin{tad}

\tadNom{FileDePriorité} 

\tadDependances{Booleen} 

\begin{tadOperations}{fileDePriorite}
\tadOperation{fileDePriorite}{}{\tadParams{FileDePriorite}}
\tadOperation{enfiler}{\tadDeuxParams{FileDePriorite}{Element}}{\tadParams{FileDePriorite}}
\tadOperation{defiler}{\tadParam{FileDePriorite}}{\tadParams{FileDePriorite}}
\tadOperation{longueur}{\tadParams{FiledePriorite}}{\tadParams{\naturel}}
\tadOperation{nouvelleFile}{\tadParams{FiledePriorite}}{\tadParams{FiledePriorite}}
\tadOperation{estVide}{\tadParams{FileDePriorite}}{\tadParams{\booleen}}
\tadOperation{estPresent}{\tadDeuxParams{FileDePriorite}{Element}}{\tadParams{\booleen}}
\tadOperation{obtenirElement}{\tadParams{FileDePriorite}}{\tadParams{Element}}
\end{tadOperations}

\begin{tadSemantiques}{fileDePriorite}
    \tadSemantiques{filedePriorite}{Créer une file de priotité vide}
    \tadSemantiques{enfiler}{Enfile un élément en fin de file}
    \tadSemantiques{defiler}{Défile le dernier élément de la file}
    \tadSemantiques{longueur}{Calcule la longueur de la file}
    \tadSemantiques{nouvelleFile}{Créer la nouvelle file de priorité à partir de l'actuelle}
    \tadSemantiques{estVide}{Cherche si la file de priotité est vide ou non}
    \tadSemantiques{estPresent}{Cherche si un élément donné est présent dans la file ou non}
    \tadSemantiques{obtenirElement}{Donne le dernier élément de la file}
\end{tadSemantiques}
    
\begin{tadPreconditions}
    formationNouvelArbrePondéré : non estVide(defiler(defiler(fdp)))\newline
\end{tadPreconditions}

\end{tad}
\section{TAD TableDeCodage}

\begin{tad}
\tadNom{TableDeCodage}
\tadDependances{ArbreHuffmann}

\begin{tadOperations}{obtenirTableDeCodage}
    \tadOperation{obtenirTableDeCodage}{\tadParams{ArbreHuffmann}}{\tadParams{TableDeCodage}}
\tadOperation{}{}{}
\end{tadOperations}

\begin{tadSemantiques}{}
    obtenirTableDeCodage : Renvoie la table de codage grâce à un arbre de Huffmann \newline
\end{tadSemantiques}
    
\begin{tadPreconditions}{obtenirTableDeCodage}
    obtenirTableDeCodage : non estVide(arbrebinaire)
\end{tadPreconditions}
\end{tad}
\end{document}






\begin{figure}
  \centering
  \includegraphics[width=0.8\linewidth]{Analyse_decsendante_compression.jpg}
  \caption{Analyse Descendante Compression}
  \label{fig:image2}
\end{figure}



\begin{figure}
  \centering
  \includegraphics[width=0.8\linewidth]{Analyse_decsendante_décompression.jpg}
  \caption{Analyse Descendante Décompression}
  \label{fig:image3}
\end{figure}


\newpage

\section{Conception préliminaire}

\subsection{\textbf{Arbre Huffman}}

\documentclass{article}
\usepackage{pseudocode}

\begin{document}
\section{TAD Arbre Huffman}
\begin{algorithme}
\signatureFonction{arbreHuffman}
	{stat : Statistiques}{Arbre Huffman}
	{}
\end{algorithme}

\begin{algorithme}
\signatureFonction{fusionner}
	{arbre1, arbre2 : Arbre Huffman}{Arbre Huffman}
	{non estVide(arbre1) ou non estVide(arbre2)}
\end{algorithme}

\begin{algorithme}
\signatureFonction{obtenirOccurences}
	{arbre : Arbre Huffman}{\naturel}
	{}
\end{algorithme}

\begin{algorithme}
	\signatureProcedure{fixerFilsGauche}
		{\paramEntreeSortie{racine : ArbreHuffman}, \paramEntree{arbreg : ArbreHuffman}}
		{}
\end{algorithme}
		
\begin{algorithme}
	\signatureProcedure{fixerFilsDroit}
		{\paramEntreeSortie{racine : ArbreHuffman}, \paramEntree{arbred : ArbreHuffman}}
		{}
\end{algorithme}

\begin{algorithme}
\signatureFonction{estFeuille}
	{arbre : Arbre Huffman}{\booleen}
	{}
\end{algorithme}

\begin{algorithme}
\signatureFonction{estVide}
	{arbre: Arbre Huffman}{\booleen}
	{}
\end{algorithme}

\begin{algorithme}
\signatureFonction{obtenirFilsGauche}
	{arbre : Arbre Huffman}{Arbre Huffman}
	{non estFeuille(arbre)}
\end{algorithme}

\begin{algorithme}
\signatureFonction{obtenirFilsDroit}
	{arbre : Arbre Huffman}{Arbre Huffman}
	{non estFeuille(arbre)}
\end{algorithme}

\begin{algorithme}
	\signatureFonction{estPresent}
		{arbre : Arbre Huffman, car : \caractere}{\booleen}
		{}
	\end{algorithme}
			
\end{document}



\subsection{\textbf{Code Binaire}}

\documentclass{article}
\usepackage{pseudocode}

\begin{document}
\section{TAD Code Binaire}
\begin{algorithme}
\signatureFonction{obtenirCodeBinaire}
	{Table : TableDeCodage, c : \caractere}{Code Binaire}
	{}
\end{algorithme}
\begin{algorithme}
\signatureFonction{estBinaire}
	{codeBinaire : Code Binaire}{\booleen}
	{non estVide(codeBinaire)}
\end{algorithme}

\end{document}


\subsection{\textbf{Statistiques}}

\documentclass{article}
\usepackage{pseudocode}

\begin{document}
\section{TAD Statistiques}
\begin{algorithme}
\signatureFonction{statistique}
	{c : \caractere, n : \naturel}{Statistiques}
	{}
\end{algorithme}
\begin{algorithme}
\signatureFonction{obtenirOccurences}
	{texte : \chaine}{Statistiques}
	{}
\end{algorithme}

\end{document}


\subsection{\textbf{Octet}}

\documentclass{article}
\usepackage{pseudocode}


\begin{document}
\section{TAD Octet}

\begin{algorithme}
\signatureFonction{ estOctet}{(o : Octet)}{\booleen}
{}
\end{algorithme}

\begin{algorithme}
\signatureFonction{ encoder}{(o : Octet, code : CodeBinaire)}{\chaine}
{}
\end{algorithme}

\begin{algorithme}
\signatureFonction{ decoder:}{(chaine : \chaine, code : CodeBinaire)}{Octet}
{}
\end{algorithme}

\begin{algorithme}
\signatureFonction{ obtenirOctet:}{(chaine : \chaine)}{Octet}
{}
\end{algorithme}

\end{document}


\subsection{\textbf{FileDePriorite}}

\documentclass{article}
\usepackage{pseudocode}


\begin{document}
\section{TAD FileDePriorite}

\begin{algorithme}
\signatureFonction{ filedePriorite}{}{FileDePriorite}
{}
\end{algorithme}

\begin{algorithme}
    \signatureProcedure{enfiler}{arbreH : ArbreHuffman, file : FileDePriorite }{FileDePriorite}
    {}
\end{algorithme}

\begin{algorithme}
        \signatureProcedure{defiler}{file : FileDePriorite }{FileDePriorite}
        {non(estVide(file))}
\end{algorithme}

\begin{algorithme}
    \signatureFonction{obtenirElement}{file : FileDePriorite }{ArbreHuffman}
    {non(estVide(file))}
\end{algorithme}

\begin{algorithme}
    \signatureFonction{obtenirArbre}{file : FileDePriorite }{ArbreHuffmann}
    {non(estVide(file))}
\end{algorithme}

\begin{algorithme}
\signatureProcedure{ nouvelleFile}{\paramEntreeSortie{file : FileDePriorite}}
{non estVide(defiler(defiler(fdp)))}
\end{algorithme}

\end{document}


\subsection{\textbf{TableDeCodage}}

\documentclass{article}
\usepackage{pseudocode}

\begin{document}
\section{TAD TableDeCodage}

\begin{algorithme}
\signatureFonction{obtenirTableDeCodage}
	{a : ArbreHuffman}{TableDeCodage}
	{non estVide(a)}
\end{algorithme}

\begin{algorithme}
	\signatureFonction{obtenirCode}
		{tab : TableDeCodage, car : \caractere}{CodeBinaire}
		{estPresent(tab,car)}
\end{algorithme}

\begin{algorithme}
	\signatureFonction{obtenirCarac}
		{tab : TableDeCodage, code : CodeBinaire}{\caractere}
		{}
\end{algorithme}

\begin{algorithme}
	\signatureProcedure{ajouter}
		{\paramEntreeSortie{ tab : TableDeCodage} \paramEntree{ car : \caractere, code : CodeBinaire}}
		{}
\end{algorithme}

\begin{algorithme}
	\signatureProcedure{retirer}
		{\paramEntreeSortie{ tab : TableDeCodage} \paramEntree{ car : \caractere}}
		{}
\end{algorithme}

\begin{algorithme}
	\signatureFonction{estPresent}
		{tab : TableDeCodage, car : \caractere}{\booleen}
		{}
\end{algorithme}

\end{document}



\subsection{\textbf{Compression}}

\begin{algorithme}
\signatureProcedure{COMP_ecrireBit}{\paramEntreeSortie{fichier : Fichier Binaire}\paramEntree{bit : \caractere}}
    {}
\end{algorithme}

\begin{algorithme}
\signatureProcedure{COMP_ecrireBitsDansFichier}{\paramEntreeSortie{fichier : Fichier Binaire}\paramEntree{bits : \chaine}}
    {}
\end{algorithme}

\begin{algorithme}
\signatureFonction{COMP_initSrcFile}{fileName : \chaine}{Fichier Texte}
    {}
\end{algorithme}

\begin{algorithme}
\signatureProcedure{COMP_CreerArbreEtEnfiler}{\paramEntreeSortie{fp : FileDePriorite}\paramEntree{stat : Statistiques}}
    {}
\end{algorithme}

\begin{algorithme}
\signatureFonction{COMP_FusionnerArbreFDP}{fp :  FileDePriorite}{ArbreHuffman}
    {}
\end{algorithme}

\begin{algorithme}
\signatureFonction{COMP_compresser}{fileName : \chaine}{Fichier Binaire}
    {}
\end{algorithme}


\subsection{\textbf{Decompression}}

\begin{algorithme}
\signatureFonction{fileFromFILE}{fp : FileDePriorite, caracter : \caractere, ponderation : \naturel}{FileDePriorite}
    {}
\end{algorithme}
    
\begin{algorithme}
\signatureFonction{lireBit}{fichier : Fichier Binaire}{\naturel}
    {}
\end{algorithme}
    
\begin{algorithme}
\signatureFonction{arbrePrincipal}{fp : FileDePriorite}{ArbreHuffman}
    {}
\end{algorithme}
    
\begin{algorithme}
\signatureFonction{decompresser}{fileName : Fichier Binaire}{Fichier Texte}
    {}
 \end{algorithme}

\newpage

\section{Conception détaillée}

\subsection{\textbf{Arbre Huffman}}

\section{Déclaration du type}
\begin{algorithme}

\begin{enregistrement}{ArbreHuffman}
    \champEnregistrement{statistiques}{Statistiques}
    \champEnregistrement{filsGauche}{\typePointeur{ArbreHuffman}}
    \champEnregistrement{filsDroit}{\typePointeur{ArbreHuffman}}
\end{enregistrement}
\end{algorithme}

 
\subsubsection{arbreDeHuffman}

\begin{algorithme}

    \fonction{arbreDeHuffman}{stat : Statistiques}{\typePointeur{ArbreHuffman}}{}
        {arbre : \typePointeur{ArbreHuffman}}
        {
        \sialors{non(estVide(arbre))}
        {
        \affecter{\champPointeur{arbre}{statistiques}}{stat}
        \affecter{\champPointeur{arbre}{filsGauche}}{NIL}
        \affecter{\champPointeur{arbre}{filsDroit}}{NIL}
        }
        \retourner{arbre}        
        }
\end{algorithme}

\subsubsection{fixerStatistique}

\begin{algorithme}

    \procedure{fixerStatistique}{\paramEntreeSortie{arbre : \typePointeur{ArbreHuffman}}\paramEntree{stat : Statistiques}}{}
    {}
    {
    \affecter{\champPointeur{arbre}{statistiques}}{stat}
    }


\end{algorithme}

\subsubsection{obtenirStat}

\begin{algorithme}

    \fonction{obtenirStat}{arbre : ArbreHuffman}{Statistiques}{}
    {}
    {
    \retourner{\champ{arbre}{statistiques}}
    }


\end{algorithme}

\subsubsection{obtenirPonderation}

\begin{algorithme}
    \fonction{obtenirPonderation}{arbre : ArbreHuffman}{\naturel}{}{}
    {
        \retourner{\champ{arbre}\champ{statistiques}{ponderation}}

    }
\end{algorithme}

\subsubsection{obtenirCaracter}

\begin{algorithme}
    \fonction{obtenirCaracter}{arbre : ArbreHuffman}{\caractere}{}{}
    {
        \retourner{\champ{arbre}\champ{statistiques}{caractere}}

    }
\end{algorithme}

\subsubsection{fusionner}

\begin{algorithme}

    \fonction{fusionner}{arbre1 , arbre2 : \typePointeur{ArbreHuffman}}{\typePointeur{ArbreHuffman}}{non (estVide(arbre1) et estVide(arbre2))}
    {res : \typePointeur{ArbreHuffman}, stat : Statistiques}
    {   
        \affecter{\champ{stat}{caractere}}{' '}
        \affecter{\champ{stat}{ponderation}}{obtenirPonderation(\typePointeur{arbre1})+obtenirPonderation(\typePointeur{arbre2})}
        \affecter{res}{arbreDeHuffman(stat)}
        \sialorssinon{obtenirPonderation(\typePointeur{arbre1})$<$obtenirPonderation(\typePointeur{arbre2})}
        {fixerFilsGauche(res,arbre1)\\
        fixerFilsDroit(res,arbre2)}
        {fixerFilsGauche(res,arbre2)\\
        fixerFilsDroit(res,arbre1)}
        \retourner{res}
        

    }
\end{algorithme}


\subsubsection{fixerFilsDroit}

\begin{algorithme}
    \procedure{fixerFilsDroit}{\paramEntreeSortie{arbre : \typePointeur{ArbreHuffman}} \paramEntree{droit : \typePointeur{ArbreHuffman}}}{}{}
        {
            \sialors{non(estVide(arbre))}
            {
            \affecter{\champPointeur{arbre}{filsDroit}}{droit}
            }
        }

    

\end{algorithme}

\subsubsection{fixerFilsGauche}

\begin{algorithme}
    \procedure{fixerFilsGauche}{\paramEntreeSortie{arbre : \typePointeur{ArbreHuffman}} \paramEntree{gauche : \typePointeur{ArbreHuffman}}}{}{}
        {
            \sialors{non(estVide(arbre))}{
            \affecter{\champPointeur{arbre}{filsGauche}}{gauche}}
        }
\end{algorithme}

\subsubsection{estVide}

\begin{algorithme}
    \fonction{estVide}{arbre : \typePointeur{ArbreHuffman}}{\booleen}{}{}
    {
        \retourner{((\champPointeur{arbre}\champ{statistiques}{caractere} = NULL) et (\champ{arbre}\champ{statistique}{ponderation} = NULL))}
    }
\end{algorithme}

\subsubsection{estFeuille}

\begin{algorithme}
    \fonction{estFeuille}{arbre : ArbreHuffman}{\booleen}{}{}
    {

    \retourner{(estVide(obtenirFilsGauche(@arbre)) et estVide(obtenirFilsDroit(&arbre)))}
    }
\end{algorithme}


\subsubsection{obtenirFilsGauche}

\begin{algorithme}
    \fonction{obtenirFilsGauche}{arbre : \typePointeur{ArbreHuffman}}{\typePointeur{ArbreHuffman}}{non estFeuille(arbre)}{}
    {
    \retourner{\champPointeur{arbre}{filsGauche}}
    }
\end{algorithme}

\subsubsection{obtenirFilsDroit}

\begin{algorithme}
    \fonction{obtenirFilsDroit}{arbre : \typePointeur{ArbreHuffman}}{\typePointeur{ArbreHuffman}}{non estFeuille(arbre)}{}
    {
    \retourner{\champPointeur{arbre}{filsDroit}}
    }
\end{algorithme}

\subsubsection{estPresent}

\begin{algorithme}
    \fonction{estPresent}{arbre : \typePointeur{ArbreHuffman}, caractere : \caractere}{\booleen}{}{}
    {   
        \sialors(non(estVide(arbre)))
        {
        \sialorssinon{\champPointeur{arbre}\champ{statistiques}{caractere} = caractere}
        {\retourner{Vrai}}
        {\retourner{estPresent(obtenirFilsDroit(arbre)) ou estPresent(obtenirFilsGauche(arbre))}}
        }
        }
\end{algorithme}


\subsection{\textbf{Code Binaire}}

\subsubsection{Déclaration du type}
\type{CodeBinaire}{\chaine}

\subsubsection{obtenirCodeBinaire}

\begin{algorithme}
    \fonction{obtenirCodeBinaire}{a : Arbre Huffman, car : \caractere}{CodeBinaire}
    {non(estVide(a))}
    {}
    {\sialorssinon{estFeuille(a)}
    {\retourner('0')}
    \retourner{ObtenirCodeBinaireR(a,car)}}

\subsubsection{obtenirCodeBinaireR}

\begin{algorithme}
    \fonction{obtenirCodeBinaireR}{a : Arbre Huffman, car : \caractere}{CodeBinaire}
    {non(estVide(a))}
    {}
    {\sialorssinon{estFeuille(a)}
    {\retourner{''}}
    {\sialorssinon{estPresent(car, obtenirFilsGauche(a))}
    {\retourner{'0' + obtenirCodeBinaireR(obtenirFilsGauche(a),car)}}
    {\retourner{'1' + obtenirCodeBinaireR(obtenirFilsDroit(a),car)}}}}
\end{algorithme}

\subsubsection{estVide}

\begin{algorithme}
    \fonction{estVide}{code : CodeBinaire}{\booleen}
    {}
    {}
    {\retourner{code = ''}}
\end{algorithme}

\subsubsection{obtenirIemeBit}

\begin{algorithme}
    \fonction{obtenirIemeBit}{code : CodeBinaire, i : \naturelNonNul}{Bit}
    {i $>$ 0 et i \leq longueur(code)}
    {}
    {\retourner{code[i]}}
\end{algorithme}

\subsubsection{longueur}

\begin{algorithme}
    \fonction{longueur}{code : CodeBinaire}{\naturel}
    {}
    {l : naturel}
    {\affecter{l}{0}
    \pourChaque{car}{code}
    {\affecter{l}{l+1}}
    \retourner{l}}
\end{algorithme}

\subsection{\textbf{Statistiques}}

\subsubsection{Déclaration du type}

\begin{algorithme}
    \begin{enregistrement}{Statistiques}
        \champEnregistrement{caractere}{\caractere}
        \champEnregistrement{ponderation}{\naturel}
    \end{enregistrement}    
\end{algorithme}

\subsubsection{statistique}

\begin{algorithme}
    \fonction{statistique}{c : \caractere, n : \naturel}{Statistiques}{}
        {stat : Statistiques}
        {\affecter{stat.caracteres}{c}
        \affecter{stat.ponderation}{n}
        \retourner{stat}}
\end{algorithme}

\subsubsection{obtenirPonderation}

\begin{algorithme}
    \fonction{obtenirPonderation}{stat : Statistiques}{\naturel}{}
        {}
        {\retourner{stat.ponderation}}
\end{algorithme}

\subsubsection{obtenirCaractere}

\begin{algorithme}
    \fonction{obtenirCaractere}{stat : Statistiques}{\caractere}{}
        {}
        {\retourner{stat.caractere}}
\end{algorithme}

\subsection{\textbf{Octet}}

\subsubsection{Déclaration du type}
\begin{algorithme}
\begin{enregistrement}{Octet}
    \champEnregistrement{Bit}{\tableau{1..8}{de}{Bit}}
\end{enregistrement}    
\end{algorithme}

\section{obtenirBitRestants}
\begin{algorithme}



\end{algorithme}

\section{incrementStartPosition}
\begin{algorithme}



\end{algorithme}

\section{obtenirStartPosition}
\begin{algorithme}



\end{algorithme}

\subsubsection{creerOctet}

\begin{algorithme}
    \fonction{creerOctet}{}{Octet}
    {}
    {res : Octet,
    i : \naturel}
    {\affecter{startPosition}{0}
    \pour{i}{1}{8}{}
        {\affecter{\champ{res}{bits}[i]}{BITA0}}
    \retourner{res}}
\end{algorithme}

\section{fixerBit}

\begin{algorithme}
    \procedure{fixerBit}{\paramEntreeSortie{octet : \typePointeur{Octet}}
                        \paramEntree{pos : \entier, valeur : Bit}}{}
    {position$leq$ 9}{}
    {
    \sialors{position $leq$ 9}
    {
    \affecter{\champPointeur{octet}{bits}[position]}
    }
    }
\end{algorithme}

\section{iemeBit}

\begin{algorithme}
    \fonction{iemeBit}{octet : Octet, i : \entier}{Bit}
    {}{}
    {
    \retourner{\champ{octet}{bits}[i]}
    }
\end{algorithme}

\section{decaler}

\begin{algorithme}
    \procedure{decaler}{\paramEntreeSortie{octet : \typePointeur{Octet}}}{}
    {}
    {i : \naturel}
    {\pour{i}{1}{8}{}
        {\affecter{\champPointeur{octet}{bits}[i]}{bits[i+1]}}
    \affecter{\champPointeur{octet}{bits}[8]}{BITA0}}
\end{algorithme}

\section{bitsToStr}

\begin{algorithme}
    \fonction{bitsToStr}{octet : \typePointeur{Octet}}{\chaine}
    {}
    {str : \typePointeur{\chaine},
    i : \naturel}
    {
    \allouer{str}
    \pour{i}{1}{8}{}
    \sialorssinon{\champPointeur{octet}{bits}[0] == BITA0}
    {\affecter{str[i]}{'0'}}
    {\affecter{str[i]}{'1'}}
    \affecter{str[8]}{'0'}
    }
\end{algorithme}

\section{strToBits}

\begin{algorithme}
    \fonction{strToBits}{chaine : \chaine}{Octet}
    {}
    {octet : Octet,
    i : \naturel,
    longueurChaine : \entier}
    {
    \affecter{octet}{creerOctet()}
    \affecter{longueurChaine}{taille(chaine)}
    \pour{i}{1}{longueurChaine-1}{}
    \sialorssinon{\champPointeur{chaine[i ]== '0'}
    {\affecter{\champ{octet}{bits}[i]}{BITA0}
    incrementStartPosition()}
    {\affecter{\champ{octet}{bits}[i]}{BITA1}
    incrementStartPosition()}
    }
\end{algorithme}

\subsubsection{encoder}

\begin{algorithme}
    \procedure{encoder}{\paramEntreeSortie{o : Octet}\paramEntree{code : CodeBinaire, i : \naturel}}
    {non estVide(code), 0 \leq i $<=$ 7}
    {j : \naturel}
    {\pour{j}{i}{i + longueur(code)}{}
    {modifier(o,j,obtenirIemeBit(code,j))}}
\end{algorithme}

\section{estVide}

\begin{algorithme}
    \fonction{estVide}{\booleen}{octet : \champPointeur{Octet}}
    {}
    {i : \naturel}
    {\pour{i}{1}{8}{}
    {\sialors{\champPointeur{octet}{bits}[i] $neq$ BITA0}
    {\retourner{Faux}}
    }
    \retourner{Vrai}
    }
\end{algorithme}



\subsection{\textbf{FileDePriorite}}

\documentclass{article}
\usepackage{pseudocode}

\begin{document}

\begin{algorithme}
    \type{ListeChaineeArbre}
        \typePointeur{Noeud}
    
    \end{algorithme}
    
    \begin{algorithme}
    \begin{enregistrement}{Noeud}
        \champEnregistrement{arbre}{ArbreHuffman}
        \champEnregistrement{ListeSuivante}{ListeChainee}
    \end{enregistrement}
\end{algorithme}

\begin{algorithme}
        \begin{enregistrement}{FileDePriorite}
            \champEnregistrement{debut}{ListeChaineeArbre}
            \champEnregistrement{fin}{ListeChaineeArbre}
            \champEnregistrement{longueur}{\naturel}
        \end{enregistrement}
        \end{algorithme}

\begin{algorithme}

    \fonction{fileDePriorite}{a : ArbreHuffman}{FileDePriorite}{}
        {file:FileDePriorite}
        {
        \affecter{\champ{file}\champ{debut}{arbre}}{a}
        \affecter{\champ{file}\champ{debut}{ListeSuivante}}{NIL}
        \affecter{\champ{file}{fin}}{NIL}
        \affecter{\champ{file}{longueur}{1}}
        }
\end{algorithme}

\begin{algorithme}

    \procedure{enfiler}{\paramEntreeSortie{file : FileDePriorite} \paramEntree{arbre : ArbreHuffman}}{}{}
        {temp1 : FileDePriorite, temp2 : FileDePriorite}
        {
        \affecter{temp1}{\champ{file}{debut}}
        \affecter{temp2}{\champ{file}{fin}}
        \tantque{obtenirOccurences(arbre)$>$obtenirOccurences(temp2)}{
        \affecter{temp1}{obtenirListeSuivante(temp1)} 
        \affecter{temp2}{obtenirListeSuivante(temp2)}}
        ajouter(temp1,arbre)
        fixerListeSuivante(temp1,temp2)
        \tantque{non(estVide(obtenirListeSuivante(temp1)))}
        {\affecter{temp1}{obtenirListeSuivante{temp1}}}
        \affecter{\champ{file}{fin}}{temp1}
        \affecter{\champ{file}{longueur}}{\champ{file}{longueur} + 1}

        }
        
\end{algorithme}



\begin{algorithme}
    \procedure{defiler}{\paramEntreeSortie{file : FileDePriorite}}{}{}
        {temp : FileDePriorite}
        {
        
        \affecter{temp}{\champ{file}{fin}}
        
        \affecter{\champ{file}{fin}}{obtenirListeSuivante(\champ{file}{fin})}      
        \affecter{\champ{file}{longueur}}{\champ{file}{longueur} - 1}
        \motclefDesallouer(temp)
        }
\end{algorithme}


\begin{algorithme}
    \fonction{obtenirElement}{file : FileDePriorite, i :\naturelNonNul}{Statistiques}{i$<$nb}{}
    {
        temp = file.debut;
        \pour{j}{1}{i-1}{}
            {temp = temp.ListeSuivante}
        \retourner{\champ{temp}\champ{arbre}{statistiques}}
    }
\end{algorithme}

\begin{algorithme}

\procedure{nouvelleFile}{\paramEntreeSortie{file : FileDePriorite}}{}{}
{
    enfiler(file,ArbreHuffman.fusionner(\champ{file}\champ{fin}{arbre},obtenirListeSuivante(\champ{file}\champ{fin}{arbre})))
    defiler(file)
    defiler(file)
}


\end{algorithme}



\end{document}

\subsection{\textbf{TableDeCodage}}

\documentclass{article}
\usepackage{pseudocode}

\begin{document}


\begin{algorithme}
\type{TableDeCodage}{\typePointeur{Noeud}}
\begin{enregistrement}{Noeud}
    \champEnregistrement{code}{CodeBinaire}
    \champEnregistrement{car}{\caractere}
    \champEnregistrement{listeSuivante}{TableDeCodage}

\end{enregistrement}
\end{algorithme}


\section{TableDeCodage}

\begin{algorithme}
    \fonction{obtenirTableDeCodage}{a : Arbre Huffman}{TableDeCodage}
    {non estVide(a)}
    {res : TableDeCodage}
    
    début
        res.car < obtenirCaractere(a.statistiques)
        res.code < obtenirCodeBinaire(res.car)
        listeSuivante < NIL
        retourner res
    fin

    {\affecter{file}{creerFileDePriorite(fic)}
    \affecter{arbreHuffman}{creerArbreHuffman(file)}
    \affecter{table}{creerTableDeCodage(arbreHuffman)}}
\end{algorithme}



\begin{algorithme}
    \fonction{obtenirCode}{tab: TableDeCodage, car : \caractere}{CodeBinaire}
    {estPresent(tab,car)}
    {}
    début
        si non(estVide(tab)) alors 
            si tab.car=car alors
                retourner tab.code
            sinon
                retourner obtenirCode(tab.listeSuivante)
            finsi
        finsi
    fin
    {}
\end{algorithme}

\begin{algorithme}
    \fonction{obtenirCarac}{tab: TableDeCodage, code : CodeBinaire}{\caractere}
    {}
    {}
    début
        si non(estVide(tab)) alors 
            si tab.code=code alors
                retourner tab.car
            sinon
                retourner obtenirCode(tab.listeSuivante)
            finsi
        finsi
    fin
    {}
\end{algorithme}

\begin{algorithme}
    \procedure{ajouter}{\paramEntreeSortie tab : TableDeCodage, \paramEntree car : \caractere, code : CodeBinaire}
    {}
    {temp : TableDeCodage}
    début
        temp < tab
        allouer(tab)
        tab.car < car
        tab.code < code
        tab.listeSuivante < temp
    fin
    {}
\end{algorithme}

\begin{algorithme}
    \procedure{supprimer}{\paramEntreeSortie tab : TableDeCodage, \paramEntree car : \caractere, code : CodeBinaire}
    {}
    {}

    {}
\end{algorithme}

\begin{algorithme}
    \procedure{estPresent}{\paramEntreeSortie tab : TableDeCodage, \paramEntree car : \caractere, code : CodeBinaire}
    {non(estVide(tab))}
    {}

    {}
\end{algorithme}

\end{document}


\subsection{\textbf{Compression}}

\subsubsection[Compresser]
\begin{algorithme}
    \fonction{compresser}{fic : FichierTexte}{fichierBinaire}
    {}
    {file : FileDePriorite;temp, table : TableDeCodage;  arbreHuffman : ArbreHuffman; fc : FichierBinaire}
    {
    \affecter{file}{creerFileDePriorite(fic)}
    \affecter{table}{creerTableDeCodage(FileDePriorite.obtenirElement(file))}
    \tantque{non(EOF(fic))}{
        \affecter{temp}{table}
        ecrireOctet(fc,encoder(obtenirCodeBinaire(temp,lireCaractere(fic))))
    }
    \retourner{fic}
    }
    

    
    
    
\end{algorithme}

\subsubsection{ecrireBit}

\begin{algorithme}
    \fonction{COMP_ecrireBit}{fic :  FichierBinaire, \caractere bit}{}
    {}
    {octet : FileDePriorite; ensC : Ensemble de \caractere; ensS : Ensemble de Statistiques}
    {
    \affecter{ensC}{obtenirCaracteres(fic)}
    \affecter{ensS}{creerStatistique(ensC)}
    \affecter{file}{filedePriorite()}
    \pourChaque{Stats}{ensS}{
        enfiler(file, arbrehuffman(stats))
    }
    \tantque{FileDePriorite.obtenirLongueur(file)$>$1}{  %OucreerFileDePriorite faire avec longueur >1 de FileDePrio
        nouvelleFile(file)}
 
    \retourner{file}
    }
\end{algorithme}

\subsubsection{ecrireBitsDansFichier}

\begin{algorithme}
    \fonction{ cecrireBitsDansFichier}{fic :  FichierBinaire, bit :Ensemble de \caractere }{TableDeCodage}{}{}{ longueur:\naturel }{
        \affecter{longueur}{0}

    }
    
    {\retourner{obtenirTableDeCodage(arbre)}}
\end{algorithme}

\subsubsection{initFichier}

\begin{algorithme}
    \fonction{ creerStatistique}{fic : FichierTexte;ensC : Ensemble de \caractere}{Ensemble de Statistiques}
    {}
    {ensStats : Ensemble de Statistiques, nbOcc : \naturel}   				
    {\pourChaque{car}{ensC}{
        \tantque{non(EOF(fic))}{
        lire(fic,c)
        \newline
        
        \sialors{c=car}{
            \affecter{nbOcc}{nbOcc + 1}

        
        Ensemble.ajouter(ensStats,statistique(car, nbOcc))}}}
    \retourner{ensStats}}
\end{algorithme}


\subsubsection{CreerArbreEtEnfiler}

\begin{algorithme}
    \fonction{obtenirCaracteres}{fic : FichierTexte}{Ensemble de \caractere}
    {}
    {ens : Ensemble de \caractere}
    {\affecter{ens,ensemble()}
    ouvrir(fic)
    \tantque{non(eof(fic))}
        {\lire(fic,car)
        \sialorssinon{not(estPresent(ens,car))}
            {ajouter(ens,car)}{}}
    \retourner{ens}}
\end{algorithme}


\subsubsection{FusionnerArbreFDP}

\begin{algorithme}
    \fonction{obtenirOccurences}{c : \caractere}{\naturel}
    {}
    {}
    {\retourner{Statistiques.obtenirOccurence(car)}}
\end{algorithme}


\subsection{\textbf{Decompression}}

\documentclass{article}
\usepackage{pseudocode}

\begin{document}

\section{Décompression}

\begin{algorithme}
    \signatureFonction{ decompresser}{fic : FichierHuffman}{FichierTexte}\newLine
    {stats_fic : Ensemble de Statistiques, file : FileDePriorite, temp : ArbreHuffman, digit,car : Caractere}
{
    \affecter{stats_fic}{obtenirStatistiques(fic)}
    \affecter{file}{obtenirFileDePriorite(stats_fic)}
    \affecter{temp}{obtenirElement(file)}
    ouvrir(fic)
    \tantque{non(eof(fic))}
        ouvrir(fic)
        \tantque{non(estUneFeuille(temp))}    
            \affecter{digit}{\lire(fic,car)}
            \sialorssinon{digit==0}
                {\affecter{temp}{obtenirFilsGauche(arbre)}}
                {\affecter{temp}{obtenirFilsDroit(arbre)}}
        fermer(fic)
        ouvrir(fictexte)
        \ecrire(fictexte, obtenirCaractere(temp))  
        fermer(fictexte)  
    fintantque
fin
}


\end{algorithme}

\section{obtenirFileDePriorite}

\begin{algorithme}
    \fonction{ obtenirFileDePriorite}{stats : Ensemble de Statistiques}{FileDePriorite}
    {}
    {ens : Ensemble de Arbre de Huffman, file : FileDePriorite, s : Statistiques, a : ArbreHuffman}
    {
        \affecter{file}{filedePriorite()}
        \pourChaque{s}{stats}
            ajouter(ens, arbrehuffman(s))
        \pourChaque{a}{ens}
            enfiler(file, a)
        \tantque{longueur>1}
            nouvelleFile(file)
        \retourner{file}
    }
\end{algorithme}

\section{obtenirStatistiques}

\begin{algorithme}
    \fonction{obtenirStatistiques}{fic :  FichierHuffman}{Ensemble de Statistiques}
    {}
    {ens : Ensemble de Statistiques}
    {ouvrir(fic)
    \tantque{str!="/"}
        {\lire{fic,car}
        \lire{fic,int}
        ajouter(ens,Statistiques.statistique(car,int))}
    \retourner{ens}
    }
\end{algorithme}

\section{creerArbreHuffman}

\begin{algorithme}
    \fonction{creerArbreHuffman}{stats : Statistiques}{Arbre Huffman}
    {}
    {}
    {\retourner{ArbreHuffman.arbrehuffman(stats)}}

\end{algorithme}

\section{insererArbre}

\begin{algorithme}
    \fonction{insererArbre}{file : FileDePriorite, arbre : Arbre Huffman}{FileDePriorite}
    {}
    {}
    {FileDePriorité.enfiler(arbre,file)
    \retourner{file}}
\end{algorithme}

\section{creerArbre}
\begin{algorithme}
    \fonction{creerArbre}{stats : Statistiques}{Arbre Huffman}
    {}
    {}
    {\retourner{Huffman.arbrehuffman(stats)}}
\end{algorithme}

\end{document}

\newpage

\section{Code C et Tests Unitaires}

\subsection{Code}
\subsubsection{\textbf{Arbre Huffman}}
\lstinputlisting[language=C, caption={ArbreHuffman.c}]{src/ArbreHuffman.c}

\subsubsection{\textbf{Code Binaire}}
\lstinputlisting[language=C, caption={CodeBinaire.c}]{src/CodeBinaire.c}

\subsubsection{\textbf{Statistiques}}
\lstinputlisting[language=C, caption={Statistiques.c}]{src/Statistiques.c}

\subsubsection{\textbf{Octet}}
\lstinputlisting[language=C, caption={Octet.c}]{src/Octet.c}

\subsubsection{\textbf{FileDePriorite}}
\lstinputlisting[language=C, caption={FileDePriorite.c}]{src/FileDePriorite.c}

\subsubsection{\textbf{TableDeCodage}}
\lstinputlisting[language=C, caption={TableDeCodage.c}]src/{TableDeCodage.c}

\subsubsection{\textbf{Compression}}
\lstinputlisting[language=C, caption={Compression.c}]{src/Compression.c}

\subsubsection{\textbf{Decompression}}
\lstinputlisting[language=C, caption={Decompression.c}]{src/Decompression.c}

\subsection{Tests unitaires}

\subsubsection{\textbf{Arbre Huffman}}
\lstinputlisting[language=C, caption={testArbreHuffman.c}]{src/tests/testArbreHuffman.c}

\subsubsection{\textbf{Code Binaire}}
\lstinputlisting[language=C, caption={testCodeBinaire.c}]{src/tests/testCodeBinaire.c}

\subsubsection{\textbf{Statistiques}}
\lstinputlisting[language=C, caption={testFileDePriorite.c}]{src/tests/testFileDePriorite.c}

\subsubsection{\textbf{Octet}}
\lstinputlisting[language=C, caption={testOctet.c}]{src/tests/testOctet.c}

\subsubsection{\textbf{FileDePriorite}}
\lstinputlisting[language=C, caption={testStatistiques.c}]{src/tests/testStatistiques.c}

\subsubsection{\textbf{TableDeCodage}}
\lstinputlisting[language=C, caption={testTableDeCodage.c}]{src/tests/testTableDeCodage.c}



\newpage


\section{Conclusion}


Pour conclure ce projet,  A FAIRE

\subsection{Répartition des tâches}

\subsubsection{Légende des Croix}

\begin{itemize}
  \item \textcolor{blue}{\ding{52}} Thomas : Croix vertes 
  \item \textcolor{yellow}{\ding{52}} Woody : Croix jaunes
  \item \textcolor{red}{\ding{52}} Achraf : Croix jaunes
  \item \ \textcolor{green}{\ding{52}} Elouen : Croix jaunes
\end{itemize}
%Dans le tabkeau on utilisera des couleurs pour savoir qui a fait quoi, prenez juste la bonne et remplacez les croix du tableau d'en dessous, on s'arrangera pour que ça rende bien dans tous les cas
%ex : pour Thomas : \textcolor{blue}{\ding{52}}


\begin{table}[ht]
  \resizebox{\textwidth}{!}{
  \begin{tabular}{|c|c|c|c|c|c|}
  \hline
& TAD & CP & CD & Développement & Tests unitaires  \\
  \hline
ArbreHuffman & \textcolor{red}{X} & \textcolor{green}{X} & \textcolor{yellow}{X} & \textcolor{yellow}{X} & \textcolor{red}{X}\\
  \hline
CodeBinaire & \textcolor{blue}{X} & \textcolor{green}{X} & \textcolor{green}{X} & \textcolor{red}{X} & \textcolor{green}{X} \\
  \hline
FileDePriorite & \textcolor{green}{X} & \textcolor{blue}{X} & \textcolor{green}{X} & \textcolor{yellow}{X} & \textcolor{blue}{X} \\
  \hline
Octet & \textcolor{red}{X} & \textcolor{blue}{X} & \textcolor{yellow}{X} & \textcolor{green}{X} & \textcolor{red}{X} \\
  \hline
Statistiques & \textcolor{blue}{X} & \textcolor{green}{X} & \textcolor{red}{X} & \textcolor{red}{X} & \textcolor{green}{X}\\
  \hline
TableDeCodage & \textcolor{green}{X} & \textcolor{yellow}{X} & \textcolor{blue}{X} & \textcolor{blue}{X} & \textcolor{yellow}{X} \\
  \hline
Compression & \textcolor{yellow}{X} & \textcolor{blue}{X} & \textcolor{blue}{X} & \textcolor{red}{X} & \\
  \hline
Decompression & \textcolor{yellow}{X} & \textcolor{red}{X} & \textcolor{green}{X} & \textcolor{blue}{X} & \\
  \hline
\end{tabular}
}
\end{table}

\end{document}
